%=====================================================================
\section{Glossary}\label{app:glossary}
%=====================================================================
Terms in alphabetical order. The section where each is first used in detail is given in brackets.

\small
\begin{longtable}{p{0.24\linewidth}p{0.70\linewidth}}
\toprule
Term & Meaning \\
\midrule
\endhead
Actuator & A device that acts on the system; here the four motor--propeller units. \\
Affine & Linear plus a constant term, $A\vect x + \vect d$. \\
Air-relative thrust & Propeller thrust computed from the propeller's speed relative to the air ($\Omega_i + r$) rather than relative to the body; matters when the body spins fast (\ref{sec:val-ijrr}). \\
Algebraic Riccati equation & The matrix equation whose solution $P$ gives the optimal LQR gain; ``discrete'' version (DARE) for sampled systems (\ref{sec:background}, \ref{sec:mpc}). \\
Attainable moment set & The set of torques the remaining rotors can produce within their thrust limits. \\
Body frame & Coordinate frame fixed to the drone (\ref{sec:model}). \\
Body rates $(p,q,r)$ & Rotation rates about the body $x$, $y$, $z$ axes: roll rate, pitch rate, yaw rate. \\
Brent's method & A reliable one-dimensional root finder that needs a bracket where the function changes sign (\ref{sec:model}). \\
Bryson's rule & Choosing LQR weights as $1/(\text{largest acceptable value})^2$ (\ref{sec:linear}). \\
Capture & Bringing the drone into the relaxed hover after a failure; a ``capture controller'' does this (\ref{sec:delay}). \\
Capture set & A set of states from which the capture controller is proven to recover (\ref{sec:certificate}). \\
Central differences & Numerical derivative $(f(x+\epsilon) - f(x-\epsilon))/2\epsilon$ (\ref{sec:background}). \\
Certificate & A computer-checked mathematical proof that a property (here, recovery) holds for a whole set of states (\ref{sec:certificate}). \\
Closed loop & The system together with its feedback controller. \\
Controllability & Whether the inputs can steer the state anywhere; tested by the rank of $[B, AB, \dots]$ or the PBH test (\ref{sec:linear}). \\
Convex problem & An optimization problem in which every local minimum is a global one, so solvers find the best solution reliably; QPs with positive semidefinite cost and LMI problems are convex. With a strictly convex cost the solution is unique. \\
CVXPY & A Python library for writing convex optimization problems; it passes them to solvers such as OSQP or Clarabel. \\
Damping ratio $\zeta$ & How quickly an oscillation dies out; 0 means never, 1 means no overshoot. \\
Delay margin & The largest pure input delay a feedback loop tolerates before becoming unstable (\ref{sec:linear}). \\
Detection delay $\tau$ & Time between the rotor failure and the moment the flight computer knows about it (\ref{sec:delay}). \\
Dictionary & The list of functions (observables) used to lift the state in EDMD (\ref{sec:koopman}). \\
EDMD / EDMDc & Extended dynamic mode decomposition (with control): fitting a linear model on a dictionary of functions of the state by least squares (\ref{sec:background}). \\
Eigenvalue & A number $\lambda$ with $A\vect v = \lambda\vect v$; describes one mode of a linear system (\ref{sec:background}). \\
Equilibrium (trim) & A state and input at which nothing changes (\ref{sec:background}). \\
Euler angles & Roll $\phi$, pitch $\theta$, yaw $\psi$: three angles that describe orientation; break down at some orientations (gimbal lock). \\
Excitation & Random changes added to the inputs during data collection so that the data show how the system responds (\ref{sec:koopman}). \\
FDI / FDD & Fault detection and isolation / diagnosis: algorithms that notice and locate a failure. \\
Feedback linearization & A nonlinear control method that cancels the nonlinear terms by the control input so the remaining loop is linear. \\
Floquet theory & Stability theory for linear systems whose coefficients repeat periodically in time (\ref{sec:model}). \\
FTC & Fault-tolerant control: control that keeps working after a failure. Active FTC switches controller after detection; passive FTC does not need detection. \\
Gramian & A matrix measuring how strongly each direction of the state can be reached (controllability) or seen (observability) in a given time (\ref{sec:linear}). \\
Gyroscopic effect & The torque needed to change the direction of an angular momentum; makes spinning bodies precess (\ref{sec:model}). \\
$H_\infty$ loop shaping & A robust linear control design that shapes the loop's frequency response and maximizes a robustness margin. \\
Healthy controller & The normal controller used with four working rotors. \\
Horizon & How far ahead MPC predicts (here 20 steps, 100\,ms). \\
INDI & Incremental nonlinear dynamic inversion: a nonlinear control method that uses measured accelerations. \\
Inertia matrix & How hard it is to change the rotation of a body about each axis. \\
Interior-point method & A family of algorithms for convex optimization that moves through the inside of the feasible set; Clarabel is one. \\
Invariant set & A set that no trajectory starting inside it can leave (\ref{sec:certificate}). \\
Jacobian & Matrix of partial derivatives; the linearization of a nonlinear function (\ref{sec:background}). \\
Koopman operator & A linear operator that advances functions of the state in time; lets nonlinear systems be described linearly in a lifted space (\ref{sec:background}). \\
Lifted state & The vector of dictionary functions $\vect z = \psi(\vect x)$. \\
LMI & Linear matrix inequality: a convex constraint saying a matrix, linear in the unknowns, is positive or negative semidefinite (\ref{sec:certificate}). \\
LPV & Linear parameter-varying: a linear model whose matrices depend on a measured, changing parameter (here the yaw rate). \\
LQR & Linear-quadratic regulator: the state-feedback gain minimizing a quadratic cost (\ref{sec:background}). \\
LTI & Linear time-invariant: a model $\dot{\vect x} = A\vect x + B\vect u$ with constant matrices. \\
Lyapunov function & An energy-like function that decreases along the motion; its existence proves stability (\ref{sec:certificate}). \\
Measured disturbance & A known, additive term in an MPC prediction model; used here for the learned residual (\ref{sec:mpc}). \\
Minkowski sum $\oplus$ & The set of all sums $a + b$ with $a$ from one set and $b$ from another (\ref{sec:certificate}). \\
Monte Carlo & Repeating a simulation from many random starts to estimate probabilities (\ref{sec:delay}). \\
Motor lag & Delay of the motor thrust behind its command, modelled as a first-order lag with time constant $\sigma$ (\ref{sec:model}). \\
MPC & Model predictive control: repeatedly optimize the inputs over a short future horizon and apply the first one (\ref{sec:background}, \ref{sec:mpc}). \\
NMPC & Nonlinear MPC: MPC with a nonlinear prediction model (no longer a QP). \\
Normalized RMSE & Root mean square error after dividing each state by its typical size (\ref{sec:koopman}). \\
Observability & Whether the state can be worked out from the sensor outputs (\ref{sec:linear}). \\
Offset-free MPC & MPC that estimates a constant disturbance and compensates it, removing steady errors like integral action. \\
Open loop & Without feedback; an open-loop prediction uses no later measurements. \\
Operator splitting & A way of solving an optimization problem by alternating between simpler subproblems; used by OSQP. \\
OSQP & A fast solver for quadratic programs. \\
Overfitting & Fitting the noise or quirks of the training data instead of the true pattern. \\
PBH test & Popov--Belevitch--Hautus test: checks controllability (or observability) mode by mode (\ref{sec:linear}). \\
Plant & The system being controlled (here the drone, or a truth plant in simulation). \\
Positive definite & A symmetric matrix $P$ with $\vect x^\top P\vect x>0$ for all $\vect x\ne0$, written $P\succ0$; $\succeq$ allows equality (\ref{sec:background}). \\
QP & Quadratic program: minimize a quadratic cost subject to linear constraints; convex. \\
Quaternion & Four numbers that describe an orientation without gimbal lock (\ref{sec:model}). \\
Randomized plants & Training on many plants with randomly drawn parameters, also called domain randomization (\ref{sec:koopman}). \\
Reaction torque & Twisting torque on the body from a spinning propeller (\ref{sec:intro}). \\
Reduced attitude $\vect n$ & The inertial ``up'' direction seen in body axes, $R^\top\vect e_3$; all that is controlled after the failure (\ref{sec:model}). \\
Relaxed hover & The spinning equilibrium of a multicopter that has lost rotors (\ref{sec:model}). \\
Residual & The difference between the truth and the physics model's prediction (\ref{sec:koopman}). \\
Ridge regression & Least squares with a penalty on the size of the coefficients (\ref{sec:background}). \\
RK4 & The classical fourth-order Runge--Kutta integration method (\ref{sec:background}). \\
Rotation matrix $R$ & Matrix turning body-frame vectors into inertial-frame vectors. \\
Saturation & Clipping of an input at its physical limits (\ref{sec:background}). \\
Schur complement & A rule that turns certain nonlinear matrix inequalities into LMIs (\ref{sec:certificate}). \\
SE(3) & The set of all positions and orientations of a rigid body. \\
Singular values & Numbers that measure how much a matrix stretches vectors in each direction; a value near zero means a direction the matrix almost loses. \\
Sliding-mode control & A robust nonlinear control method that forces the state onto a chosen surface by switching control. \\
Spectral radius & The largest eigenvalue magnitude of a discrete-time matrix; above 1 means growth. \\
Spin-up & Flying on two opposing rotors with no attitude feedback until the drone spins fast enough for the failure LQR (\ref{sec:delay}). \\
S-procedure & A standard technique that turns ``one quadratic inequality implies another'' into an LMI (\ref{sec:certificate}). \\
Standardize & Subtract the mean and divide by the standard deviation. \\
State & The smallest set of numbers that determines the future of a system given its inputs (\ref{sec:background}). \\
Sum-of-squares (SOS) programming & An extension of LMI methods that certifies polynomial inequalities, used for nonlinear stability regions. \\
Supervisor & The logic that switches from the healthy to the failure controller after detection (\ref{sec:delay}). \\
Switched system & A system that jumps between several models, for example healthy and faulty; ``dwell time'' is the minimum average time between switches for which stability is guaranteed. \\
Thrust & The upward force of a propeller, $f = \kappa_f\Omega^2$. \\
Truth plant & A simulated plant with effects the nominal model lacks, used to generate data and to test controllers (\ref{sec:truth}). \\
Wilson interval & A confidence interval for a probability estimated from a count (\ref{sec:delay}). \\
Zero-order hold (ZOH) & Holding a digital command constant until the next sample (\ref{sec:background}). \\
\bottomrule
\end{longtable}
\normalsize
